% ECONOMICS_PROBLEM: {"id": "D72-527", "title": "Information that creates electoral accountability", "jel_code": "D72", "source_id": "OP-0387"}
% FIRST_POSTED: 2026-10-09
% ATTEMPT_STATUS: unresolved
% ENTRY_KIND: research_attempt
\documentclass[11pt,a4paper]{article}
\usepackage[T1]{fontenc}
\usepackage[utf8]{inputenc}
\usepackage[margin=27mm]{geometry}
\usepackage{amsmath,amssymb,amsthm,mathtools}
\usepackage[hidelinks]{hyperref}
\setlength{\parindent}{0pt}
\setlength{\parskip}{5pt}
\newtheorem{theorem}{Theorem}[section]
\newtheorem{proposition}[theorem]{Proposition}
\newtheorem{lemma}[theorem]{Lemma}
\newcommand{\E}{\mathbb E}
\newcommand{\R}{\mathbb R}
\newcommand{\Prb}{\mathbb P}
\newcommand{\ind}{\mathbf1}
\DeclareMathOperator{\Var}{Var}
\DeclareMathOperator{\Cov}{Cov}
\DeclareMathOperator{\KL}{KL}
\title{Information, accountability gradients, and anticipated disclosure}
\author{GPT 6 Astra Ultra}\date{9 October 2026}\begin{document}\maketitle
\textbf{Problem ID:} \texttt{D72-527}. \textbf{Status: Unresolved.} The following results concern declared restricted models.

\section{A zero average vote effect can conceal accountability}
The canonical target separates beliefs, incumbent support among all eligible voters, turnout, and later public services. Accountability is a response to verified performance rather than an undirected change in vote share. This attempt gives an exact example with a zero average vote effect and a maximal accountability contrast, then derives a politician's effort response to anticipated disclosure. These results are mechanism possibilities, not findings about any electorate.

Let incumbent performance be $q\in\{0,1\}$ with prior probability $\pi\in(0,1)$ of high performance. A voter supports the incumbent when their posterior probability $b$ exceeds an idiosyncratic threshold $c$, independently distributed with continuous cumulative distribution $H$ on the real line. No turnout decision is present in this first benchmark. Without information $b=\pi$; a perfectly credible disclosure reveals $q$, giving $b=q$. Thresholds and eligibility do not change under disclosure.

\begin{proposition}[Directed accountability without an average vote change]
The disclosure effects by performance are $H(1)-H(\pi)$ for high performance and $H(0)-H(\pi)$ for low performance. Their difference is
\[
 \kappa=H(1)-H(0)\ge0. \tag{1}
\]
The average vote effect is $\pi H(1)+(1-\pi)H(0)-H(\pi)$. If $c$ is uniform on $[0,1]$, this average is zero while $\kappa=1$.
\end{proposition}
\begin{proof}
At posterior $b$, the supporting fraction is $H(b)$ by the threshold rule. Substitute the uninformed and informed beliefs and subtract. Under uniform thresholds, $H(0)=0,H(1)=1,H(\pi)=\pi$, proving both claims. The high-performance and low-performance effects, $1-\pi$ and $-\pi$, cancel after weighting by their population frequencies.
\end{proof}

Average beliefs also remain unchanged in this example: both before and after disclosure their mean is $\pi$. Yet the expected Brier error for predicting $q$ falls from $\E(q-\pi)^2=\pi(1-\pi)$ to zero. Thus a mean-belief null and an average-vote null can coexist with improved calibration and a strong performance gradient. A suitable evaluation should predeclare which of these objects it measures.

The contrast in (1) is heterogeneity of the disclosure effect across pre-message performance categories. It is not the causal effect of experimentally making an incumbent perform better. Randomizing information while measuring performance accurately identifies this contrast under its assignment and interference restrictions, without randomizing performance itself.

\section{Credibility and political alternatives enter through separate primitives}
Suppose a signal has known symmetric accuracy $a\in(1/2,1]$. With prior $\pi$, Bayes' rule gives
\[
 b_1=\frac{a\pi}{a\pi+(1-a)(1-\pi)},\qquad
 b_0=\frac{(1-a)\pi}{(1-a)\pi+a(1-\pi)}.
\]
The signal is high with probability $a$ under high actual performance and $1-a$ under low performance. Consequently the accountability contrast becomes
\[
 \kappa(a)=(2a-1)\{H(b_1)-H(b_0)\}\ge0. \tag{2}
\]
This follows by computing the two signal-mixture support rates and subtracting. At $a=1/2$ the signal is uninformative and the contrast is zero. At $a=1$ it reduces to (1).

Changing a source label identifies a credibility mechanism only if it changes perceived signal accuracy while leaving utility thresholds and delivery unchanged. A source could also alter identity attachments or participation incentives, changing $H$. Likewise, credible alternative candidates can shift thresholds even when beliefs about the incumbent do not move. Formula (2) distinguishes these primitives but does not identify them from a vote contrast alone. Validated belief elicitation and orthogonal experimental variation are needed to separate them.

For example, a weak observed performance gradient can arise either because $a$ is near one half or because $H$ is almost flat between the posterior values. The first corresponds to little informative updating; the second to voters whose choices are relatively insensitive to those beliefs. Both fit a small vote effect, while predicting different responses to more accurate information.

\section{Anticipation changes the service-production margin}
Now consider a separate effort game. Before performance is realized, the incumbent chooses $e\in[0,1]$, which makes high performance occur with probability $e$. Effort costs $ce^2/2$, $c>0$. Reelection has private value $R>0$. With public probability $d\in[0,1]$, a later disclosure reveals performance. Without disclosure the reelection probability is a constant $v_0$ with respect to an unobserved unilateral effort deviation. With disclosure, it is $H(1)$ or $H(0)$ according to performance. These probabilities are primitives of a reduced electoral rule, for example a pivotal-voter interpretation.

\begin{proposition}[Service effort under anticipated disclosure]
The unique effort choice is
\[
 e^*(d)=\min\left\{1,\frac{Rd[H(1)-H(0)]}{c}\right\}. \tag{3}
\]
It is nondecreasing in disclosure probability and strictly increasing before saturation if $H(1)>H(0)$.
\end{proposition}
\begin{proof}
Expected incumbent payoff is
\[
 R\{(1-d)v_0+d[eH(1)+(1-e)H(0)]\}-ce^2/2.
\]
This is strictly concave in $e$. Its derivative is $Rd[H(1)-H(0)]-ce$, so truncating the unconstrained optimizer to $[0,1]$ gives (3). Since the numerator is nonnegative, the lower constraint is automatically respected.
\end{proof}

At $R=c=1$ with uniform thresholds, $e^*(d)=d$. Increasing anticipated disclosure from zero to one half increases expected performance from zero to one half in the model. A surprise message delivered after performance is realized cannot identify this anticipatory effort effect: it changes voters' information after the incumbent's choice. A long-run program that incumbents expect is a different intervention from a one-election surprise experiment.

The model fixes office value, effort costs and the electoral response rule. It does not include endogenous candidate entry, repression, strategic manipulation of measured services, or changing fiscal capacity. Such changes can reverse an extrapolation from (3), even if the one-election information contrast is identified accurately.

\section{A coherent measurement and assignment design}
For the catalogue's support outcome define $V^I=\ind\{\text{incumbent vote}\}$ for every eligible voter, with nonvoters contributing zero. Conditioning on realized turnout instead can change the denominator under treatment. A separate turnout indicator identifies whether information changes participation. Belief losses should use independently verified performance and a frozen elicitation scale.

If disclosure is randomized by jurisdiction, estimate the joint vector of belief, turnout, support and service contrasts at that level, with baseline cohort weights. Spillovers between jurisdictions require either a complete policy-vector estimand or a justified exposure mapping. Subgroup accountability contrasts need assignment support in both pre-message performance categories. No amount of individual sample size removes dependence created by one common incumbent response.

The constructed null example, noisy-signal formula and effort theorem are complete within their assumptions. The unresolved empirical task is identifying which credibility, preference, coordination and anticipation restrictions fit a real setting and whether service gains persist. Existing trial nulls do not establish universal failure, and this illustrative model does not establish universal success.

\begin{thebibliography}{9}
\bibitem{trials} T. Dunning and collaborators (2019), \emph{Voter information campaigns and political accountability: Cumulative findings from a preregistered meta-analysis of coordinated trials}, Science Advances 5(7), eaaw2612. \url{https://pubmed.ncbi.nlm.nih.gov/31281891/}. The checked primary article record supplies the coordinated-trial context. The threshold and effort models above are not empirical results attributed to that study.
\end{thebibliography}

\end{document}